\chapter{2017年全国III卷（文）}

\section{单选题}

\begin{problem}
已知集合 \(A=\{1,2,3,4\}\)，\(B=\{2,4,6,8\}\)，则 \(A\cap B\) 中元素的个数为（\quad）

\choices
    {\(1\)}
    {\(2\)}
    {\(3\)}
    {\(4\)}
\end{problem}

\begin{answer}
B
\end{answer}

\begin{solution}
集合交集由同时属于两个集合的元素组成.这里 \(A\) 与 \(B\) 的公共元素为 \(2\)，\(4\)，所以 \(A\cap B=\{2,4\}\)，其中有 \(2\) 个元素，故选 B.
\end{solution}

\begin{problem}
复平面内表示复数 \(z=\i(-2+\i)\) 的点位于（\quad）

\choices
    {第一象限}
    {第二象限}
    {第三象限}
    {第四象限}
\end{problem}

\begin{answer}
C
\end{answer}

\begin{solution}
化简复数得 \(z=\i(-2+\i)=-2\i+\i^2=-1-2\i\).它在复平面中对应的点为 \((-1,-2)\)，横坐标、纵坐标均为负，因此位于第三象限，故选 C.
\end{solution}

\begin{problem}
某城市为了解游客人数的变化规律，提高旅游服务质量，收集并整理了 2014 年 1 月至 2016 年 12 月期间月接待游客量（单位：万人）的数据，绘制了下面的折线图.

\bitmapfigure[max width=0.9\linewidth]{img/2017/national_paper_3_liberal/q3_fig1.png}

根据该折线图，下列结论错误的是（\quad）

\choices
    {月接待游客量逐月增加}
    {年接待游客量逐年增加}
    {各年的月接待游客量高峰期大致在 \(7\)，\(8\) 月}
    {各年 \(1\) 月至 \(6\) 月的月接待游客量相对于 \(7\) 月至 \(12\) 月，波动性更小，变化比较平稳}
\end{problem}

\begin{answer}
A
\end{answer}

\begin{solution}
由折线图可见，月接待游客量并非逐月增加，例如 2014 年 8 月至 9 月出现下降，所以结论 A 错误.图中各年数据的总体水平逐年提高，故年接待游客量逐年增加；各年的高点大致出现在 \(7\)，\(8\) 月；与 \(7\) 月至 \(12\) 月相比，各年 \(1\) 月至 \(6\) 月的折线变化较平稳.因此错误的结论是 A.
\end{solution}

\begin{problem}
已知 \(\sin\alpha-\cos\alpha=\dfrac{4}{3}\)，则 \(\sin2\alpha=\)（\quad）

\choices
    {\(-\dfrac{7}{9}\)}
    {\(-\dfrac{2}{9}\)}
    {\(\dfrac{2}{9}\)}
    {\(\dfrac{7}{9}\)}
\end{problem}

\begin{answer}
A
\end{answer}

\begin{solution}
两边平方得 \((\sin\alpha-\cos\alpha)^2=\dfrac{16}{9}\).又 \((\sin\alpha-\cos\alpha)^2=\sin^2\alpha+\cos^2\alpha-2\sin\alpha\cos\alpha=1-\sin2\alpha\)，所以 \(1-\sin2\alpha=\dfrac{16}{9}\)，从而 \(\sin2\alpha=-\dfrac79\)，故选 A.
\end{solution}

\begin{problem}
设 \(x\)，\(y\) 满足约束条件
\examdisplaycases{}{
3x+2y-6\le 0,\\
x\ge 0,\\
y\ge 0
}
则 \(z=x-y\) 的取值范围是（\quad）

\choices
    {\([-3,0]\)}
    {\([-3,2]\)}
    {\([0,2]\)}
    {\([0,3]\)}
\end{problem}

\begin{answer}
B
\end{answer}

\begin{solution}
可行域是由 \(x\geq0\)，\(y\geq0\)，\(3x+2y\leq6\) 围成的三角形，其顶点为 \((0,0)\)，\((2,0)\)，\((0,3)\).在线性目标函数 \(z=x-y\) 的可行域上，最值出现在顶点处，分别得到 \(0\)，\(2\)，\(-3\).因此 \(z\in[-3,2]\)，故选 B.
\end{solution}

\begin{problem}
函数 \(f(x)=\dfrac15\sin\left(x+\dfrac{\pi}{3}\right)+\cos\left(x-\dfrac{\pi}{6}\right)\) 的最大值为（\quad）

\choices
    {\(\dfrac65\)}
    {\(1\)}
    {\(\dfrac35\)}
    {\(\dfrac15\)}
\end{problem}

\begin{answer}
A
\end{answer}

\begin{solution}
利用 \(\cos u=\sin\left(\dfrac\pi2+u\right)\)，有
\(\cos\left(x-\dfrac\pi6\right)=\sin\left(x+\dfrac\pi3\right)\).因此
\(f(x)=\dfrac15\sin\left(x+\dfrac\pi3\right)+\sin\left(x+\dfrac\pi3\right)=\dfrac65\sin\left(x+\dfrac\pi3\right)\).由于正弦函数的最大值为 \(1\)，所以 \(f(x)\) 的最大值为 \(\dfrac65\)，故选 A.
\end{solution}

\begin{problem}
函数 \(y=1+x+\dfrac{\sin x}{x^2}\) 的部分图象大致为（\quad）

\choices
    {\choicebitmap[width=0.15\paperwidth]{img/2017/national_paper_3_liberal/q7_optA.png}}
    {\choicebitmap[width=0.15\paperwidth]{img/2017/national_paper_3_liberal/q7_optB.png}}
    {\choicebitmap[width=0.15\paperwidth]{img/2017/national_paper_3_liberal/q7_optC.png}}
    {\choicebitmap[width=0.15\paperwidth]{img/2017/national_paper_3_liberal/q7_optD.png}}
\end{problem}

\begin{answer}
D
\end{answer}

\begin{solution}
函数的定义域为 \(x\ne0\).当 \(x\to0^-\) 时，\(\dfrac{\sin x}{x^2}\sim\dfrac1x\to-\infty\)；当 \(x\to0^+\) 时，\(\dfrac{\sin x}{x^2}\to+\infty\).当 \(x\to\pm\infty\) 时，\(\dfrac{\sin x}{x^2}\to0\)，所以两支都趋近于斜渐近线 \(y=x+1\).

再求导得
\[
f'(x)=1+\dfrac{x\cos x-2\sin x}{x^3}
\]
当 \(x\to0^+\) 时，\(f'(x)\to-\infty\)，而 \(f'(\pi)=1-\dfrac1{\pi^2}>0\)，故右支在靠近 \(x=0\) 处下降，随后出现极小值并上升.同理，\(f'(-\pi)=1-\dfrac1{\pi^2}>0\)，而 \(f'(x)\to-\infty\)（\(x\to0^-\)），故左支先上升后下降至 \(-\infty\).结合竖直渐近线两侧的极限和斜渐近线，只有选项 D 符合，故选 D.
\end{solution}

\begin{problem}
执行如图的程序框图，为使输出 \(S\) 的值小于 \(91\)，则输入的正整数 \(N\) 的最小值为（\quad）

\bitmapfigure[max width=0.62\linewidth]{img/2017/national_paper_3_liberal/q8_fig1.png}

\choices
    {\(5\)}
    {\(4\)}
    {\(3\)}
    {\(2\)}
\end{problem}

\begin{answer}
D
\end{answer}

\begin{solution}
程序开始时 \(t=1\)，\(M=100\)，\(S=0\).当 \(N=1\) 时循环一次，输出 \(S=100\)，不小于 \(91\).当 \(N=2\) 时，第一次循环后 \(S=100\)，\(M=-10\)；第二次循环后 \(S=100-10=90<91\).由于 \(N\) 为正整数，满足条件的最小值为 \(2\)，故选 D.
\end{solution}

\begin{problem}
已知圆柱的高为 \(1\)，它的两个底面的圆周在直径为 \(2\) 的同一个球的球面上，则该圆柱的体积为（\quad）

\choices
    {\(\pi\)}
    {\(\dfrac{3\pi}{4}\)}
    {\(\dfrac{\pi}{2}\)}
    {\(\dfrac{\pi}{4}\)}
\end{problem}

\begin{answer}
B
\end{answer}

\begin{solution}
球的半径为 \(1\).取圆柱的轴截面，得到高为 \(1\)、底面直径为 \(2r\) 的矩形，其外接圆半径为球半径，因此
\(r^2+\left(\dfrac12\right)^2=1\)，解得 \(r^2=\dfrac34\).所以圆柱体积为 \(V=\pi r^2\cdot1=\dfrac{3\pi}{4}\)，故选 B.
\end{solution}

\begin{problem}
在正方体 \(ABCD-A_1B_1C_1D_1\) 中，\(E\) 为棱 \(CD\) 的中点，则（\quad）

\choices
    {\(A_1E\perp DC_1\)}
    {\(A_1E\perp BD\)}
    {\(A_1E\perp BC_1\)}
    {\(A_1E\perp AC\)}
\end{problem}

\begin{answer}
C
\end{answer}

\begin{solution}
设正方体棱长为 \(a\)，取坐标
\(A(0,0,0)\)，\(B(a,0,0)\)，\(C(a,a,0)\)，\(D(0,a,0)\)，\(A_1(0,0,a)\)，\(C_1(a,a,a)\).则 \(E\left(\dfrac a2,a,0\right)\)，从而
\(\overrightarrow{A_1E}=\left(\dfrac a2,a,-a\right)\).四个选项中，
\(\overrightarrow{DC_1}=(a,0,a)\)，\(\overrightarrow{BD}=(-a,a,0)\)，\(\overrightarrow{BC_1}=(0,a,a)\)，\(\overrightarrow{AC}=(a,a,0)\).
相应的数量积依次为
\(-\dfrac{a^2}{2}\)，\(\dfrac{a^2}{2}\)，\(0\)，\(\dfrac{3a^2}{2}\).
由于 \(a>0\)，只有第三个数量积为零，所以 \(A_1E\perp BC_1\)，故选 C.
\end{solution}

\begin{problem}
已知椭圆 \(C:\dfrac{x^2}{a^2}+\dfrac{y^2}{b^2}=1\)（\(a>b>0\)）的左、右顶点分别为 \(A_1\)，\(A_2\)，且以线段 \(A_1A_2\) 为直径的圆与直线 \(bx-ay+2ab=0\) 相切，则 \(C\) 的离心率为（\quad）

\choices
    {\(\dfrac{\sqrt6}{3}\)}
    {\(\dfrac{\sqrt3}{3}\)}
    {\(\dfrac{\sqrt2}{3}\)}
    {\(\dfrac13\)}
\end{problem}

\begin{answer}
A
\end{answer}

\begin{solution}
以 \(A_1A_2\) 为直径的圆的圆心为原点，半径为 \(a\)，其方程为 \(x^2+y^2=a^2\).直线 \(bx-ay+2ab=0\) 与该圆相切，所以原点到直线的距离等于半径，即
\(\dfrac{2ab}{\sqrt{a^2+b^2}}=a\).因 \(a>0\)，可约去 \(a\)，得 \(2b=\sqrt{a^2+b^2}\)，于是 \(a^2=3b^2\).椭圆焦半距满足 \(c^2=a^2-b^2=2b^2\)，所以离心率为 \(e=\dfrac ca=\dfrac{\sqrt2}{\sqrt3}=\dfrac{\sqrt6}{3}\)，故选 A.
\end{solution}

\begin{problem}
已知函数 \(f(x)=x^2-2x+a\bigl(\e^{x-1}+\e^{1-x}\bigr)\) 有唯一零点，则 \(a=\)（\quad）

\choices
    {\(-\dfrac12\)}
    {\(\dfrac13\)}
    {\(\dfrac12\)}
    {\(1\)}
\end{problem}

\begin{answer}
C
\end{answer}

\begin{solution}
令 \(t=x-1\)，则
\(f(x)=t^2-1+a\left(\e^t+\e^{-t}\right)\).这是关于 \(t\) 的偶函数.若它有唯一零点，该零点只能是 \(t=0\)，否则非零零点 \(t\) 与 \(-t\) 会同时出现.因此 \(f(1)=0\)，即 \(-1+2a=0\)，得 \(a=\dfrac12\).反之，当 \(a=\dfrac12\) 时，\(t=0\) 是零点；当 \(t\ne0\) 时，\(t^2>0\) 且 \(\e^t+\e^{-t}>2\)，故 \(f(x)>0\)，所以零点确实唯一.故选 C.
\end{solution}

\section{填空题}

\begin{problem}
已知向量 \(\bs a=(-2,3)\)，\(\bs b=(3,m)\)，且 \(\bs a\perp\bs b\)，则 \(m=\fillinblank{}\).
\end{problem}

\begin{answer}
\(2\)
\end{answer}

\begin{solution}
向量垂直等价于数量积为零，因此
\(\bs a\cdot\bs b=(-2)\cdot3+3m=0\).解得 \(m=2\).
\end{solution}

\begin{problem}
双曲线 \(\dfrac{x^2}{a^2}-\dfrac{y^2}{9}=1\)（\(a>0\)）的一条渐近线方程为 \(y=\dfrac35x\)，则 \(a=\fillinblank{}\).
\end{problem}

\begin{answer}
\(5\)
\end{answer}

\begin{solution}
双曲线 \(\dfrac{x^2}{a^2}-\dfrac{y^2}{9}=1\) 的渐近线为 \(y=\pm\dfrac3a x\).题给渐近线斜率为 \(\dfrac35\)，所以 \(\dfrac3a=\dfrac35\)，解得 \(a=5\).
\end{solution}

\begin{problem}
\(\triangle ABC\) 的内角 \(A\)，\(B\)，\(C\) 的对边分别为 \(a\)，\(b\)，\(c\)，已知 \(C=60^\circ\)，\(b=\sqrt6\)，\(c=3\)，则 \(A=\fillinblank{}\).
\end{problem}

\begin{answer}
\(75^\circ\)
\end{answer}

\begin{solution}
由正弦定理，
\(\dfrac{b}{\sin B}=\dfrac{c}{\sin C}\)，所以
\(\sin B=\dfrac{b\sin C}{c}=\dfrac{\sqrt6\cdot\frac{\sqrt3}{2}}{3}=\dfrac{\sqrt2}{2}\).故 \(B=45^\circ\) 或 \(135^\circ\).由于 \(c=3>b=\sqrt6\)，应有 \(C>B\)，排除 \(B=135^\circ\)，得 \(B=45^\circ\).于是 \(A=180^\circ-60^\circ-45^\circ=75^\circ\).
\end{solution}

\begin{problem}
设函数
\examdisplaycases{f(x)=}{
x+1, & x\le 0,\\
2^x, & x>0
}
则满足 \(f(x)+f\left(x-\dfrac12\right)>1\) 的 \(x\) 的取值范围是\fillinblank{}.
\end{problem}

\begin{answer}
\(\left(-\dfrac14,+\infty\right)\)
\end{answer}

\begin{solution}
分三段讨论 \(x\) 与 \(x-\dfrac12\) 所在的区间.当 \(x\leq0\) 时，
\(f(x)+f\left(x-\dfrac12\right)=(x+1)+\left(x-\dfrac12+1\right)=2x+\dfrac32\)，不等式化为 \(x>-\dfrac14\)，故得到 \(\left(-\dfrac14,0\right]\).

当 \(0<x\leq\dfrac12\) 时，左边为 \(2^x+x+\dfrac12\)，因 \(2^x>1\) 且 \(x>0\)，恒大于 \(1\).当 \(x>\dfrac12\) 时，左边为 \(2^x+2^{x-\frac12}>1\)，也恒成立.合并得 \(x\in\left(-\dfrac14,+\infty\right)\).
\end{solution}

\section{解答题}

\begin{problem}
设数列 \(\{a_n\}\) 满足 \(a_1+3a_2+\cdots+(2n-1)a_n=2n\).

\begin{enumerate}
\item 求 \(\{a_n\}\) 的通项公式；
\item 求数列 \(\left\{\dfrac{a_n}{2n+1}\right\}\) 的前 \(n\) 项和.
\end{enumerate}
\end{problem}

\begin{solution}
\begin{enumerate}
\item 令题中等式在 \(n\) 与 \(n-1\) 时分别成立.两式相减，得
\( (2n-1)a_n=2n-2(n-1)=2\)，所以当 \(n\geq2\) 时，\(a_n=\dfrac{2}{2n-1}\).当 \(n=1\) 时，原等式给出 \(a_1=2=\dfrac{2}{2\cdot1-1}\)，也符合上式.因此
\(a_n=\dfrac{2}{2n-1}\)，\(n\in\mathbb N^*\).
\item 设 \(b_n=\dfrac{a_n}{2n+1}\)，则
\(b_n=\dfrac{2}{(2n-1)(2n+1)}=\dfrac{1}{2n-1}-\dfrac{1}{2n+1}\).于是前 \(n\) 项和为
\(b_1+b_2+\cdots+b_n=\left(1-\dfrac13\right)+\left(\dfrac13-\dfrac15\right)+\cdots+\left(\dfrac{1}{2n-1}-\dfrac{1}{2n+1}\right)=1-\dfrac{1}{2n+1}=\dfrac{2n}{2n+1}\).
\end{enumerate}
\end{solution}

\begin{problem}
某超市计划按月订购一种酸奶，每天进货量相同，进货成本每瓶 \(4\text{ 元}\)，售价每瓶 \(6\text{ 元}\)，未售出的酸奶降价处理，以每瓶 \(2\text{ 元}\) 的价格当天全部处理完. 根据往年销售经验，每天需求量与当天最高气温（单位：\(^{\circ}\mathrm{C}\)）有关. 如果最高气温不低于 \(25\)，需求量为 \(500\text{ 瓶}\)；如果最高气温位于区间 \([20,25)\)，需求量为 \(300\text{ 瓶}\)；如果最高气温低于 \(20\)，需求量为 \(200\text{ 瓶}\). 为了确定六月份的订购计划，统计了前三年六月份各天的最高气温数据，得下面的频数分布表：

\begin{center}
\begin{tabular}{c|cccccc}
最高气温 & \([10,15)\) & \([15,20)\) & \([20,25)\) & \([25,30)\) & \([30,35)\) & \([35,40)\) \\
\hline
天数 & \(2\) & \(16\) & \(36\) & \(25\) & \(7\) & \(4\)
\end{tabular}
\end{center}

以最高气温位于各区间的频率代替最高气温位于该区间的概率.

\begin{enumerate}
\item 估计六月份这种酸奶一天的需求量不超过 \(300\) 瓶的概率；
\item 设六月份一天销售这种酸奶的利润为 \(Y\)（单位：元），当六月份这种酸奶一天的进货量为 \(450\) 瓶时，写出 \(Y\) 的所有可能值，并估计 \(Y\) 大于零的概率.
\end{enumerate}
\end{problem}

\begin{solution}
\begin{enumerate}
\item 三年六月共有 \(2+16+36+25+7+4=90\) 天.需求量不超过 \(300\) 瓶对应最高气温低于 \(20\) 或位于 \([20,25)\)，所以对应天数为 \(2+16+36=54\)，概率估计为 \(\dfrac{54}{90}=\dfrac35\).
\item 进货量为 \(450\) 瓶，进货成本为 \(4\times450=1800\) 元.若需求量为 \(200\) 瓶，销售收入为 \(6\times200\) 元，剩余 \(250\) 瓶按每瓶 \(2\) 元处理，因此
\(Y=6\times200+2\times250-4\times450=-100\).若需求量为 \(300\) 瓶，剩余 \(150\) 瓶，故
\(Y=6\times300+2\times150-4\times450=300\).若需求量为 \(500\) 瓶，全部 \(450\) 瓶按售价售出，故 \(Y=6\times450-4\times450=900\).所以 \(Y\) 的所有可能值为 \(-100\)，\(300\)，\(900\).其中 \(Y>0\) 对应需求量为 \(300\) 或 \(500\) 瓶，对应天数为 \(36+25+7+4=72\)，故
\(P(Y>0)=\dfrac{72}{90}=\dfrac45\).
\end{enumerate}
\end{solution}

\begin{problem}
如图，四面体 \(ABCD\) 中，\(\triangle ABC\) 是正三角形，\(AD=CD\).

\begin{enumerate}
\item 证明：\(AC\perp BD\)；
\item 已知 \(\triangle ACD\) 是直角三角形，\(AB=BD\)，若 \(E\) 为棱 \(BD\) 上与 \(D\) 不重合的点，且 \(AE\perp EC\)，求四面体 \(ABCE\) 与四面体 \(ACDE\) 的体积比.
\end{enumerate}

\FigureLayoutDeclare{img/2017/national_paper_3_liberal/q19_fig1.png}{6.0cm}{0bp 0bp 0bp 0bp}
\bitmapfigure[max width=0.52\linewidth]{img/2017/national_paper_3_liberal/q19_fig1.png}
\end{problem}

\begin{solution}
\begin{enumerate}
\item 取 \(AC\) 的中点 \(F\)，连接 \(DF\)，\(BF\).由于 \(AD=CD\)，在等腰三角形 \(ACD\) 中，中线 \(DF\) 也是高，所以 \(DF\perp AC\).又因为 \(ABC\) 是正三角形，\(BF\) 是中线，也是高，所以 \(BF\perp AC\).直线 \(DF\)，\(BF\) 相交于 \(F\)，且都在平面 \(BDF\) 内，因此 \(AC\perp\text{平面 }BDF\)，从而 \(AC\perp BD\).
\item 设 \(AC=AB=BC=s\).由于 \(AD=CD\)，且三角形 \(ACD\) 是直角三角形，直角只能在 \(D\) 处，故
\(AD=CD=\dfrac{s}{\sqrt2}\).又 \(AB=BD\)，所以 \(BD=s\).由 \(AB=BC\)，\(AD=CD\) 及公共边 \(BD\)，三角形 \(ABD\) 与三角形 \(CBD\) 全等，从而对应角相等.因 \(E\) 在 \(BD\) 上，三角形 \(ADE\) 与三角形 \(CDE\) 的两边 \(AD=CD\)，\(DE\) 为公共边，且夹角相等，所以两三角形全等，得到 \(AE=CE\).结合 \(AE\perp EC\)，三角形 \(AEC\) 是等腰直角三角形，故
\(AE=CE=\dfrac{AC}{\sqrt2}=\dfrac{s}{\sqrt2}=CD\).

三角形 \(BCD\) 中 \(BC=BD=s\)，所以其底角满足 \(\angle BCD=\angle CDB\).又 \(CE=CD\)，且 \(DE\) 与 \(DB\) 同向，所以三角形 \(CDE\) 的两个底角满足 \(\angle CED=\angle CDE=\angle CDB\).因此三角形 \(CDE\sim\triangle BCD\)，且对应边给出
\(\dfrac{DE}{CD}=\dfrac{CD}{BC}=\dfrac{1}{\sqrt2}\).于是 \(DE=\dfrac{s}{2}\)，而 \(BD=s\)，故 \(BE=BD-DE=\dfrac{s}{2}=DE\)，即 \(E\) 是 \(BD\) 的中点.

以三角形 \(ACE\) 为公共底面，四面体 \(ABCE\) 与四面体 \(ACDE\) 的高分别为点 \(B\)，\(D\) 到平面 \(ACE\) 的距离.由于 \(B\)，\(E\)，\(D\) 共线且 \(E\) 在平面 \(ACE\) 内，这两个距离之比等于 \(BE:DE=1:1\)，所以体积比为 \(1\).
\end{enumerate}
\end{solution}

\begin{problem}
在直角坐标系 \(xOy\) 中，曲线 \(y=x^2+mx-2\) 与 \(x\) 轴交于 \(A\)，\(B\) 两点，点 \(C\) 的坐标为 \((0,1)\)，当 \(m\) 变化时，解答下列问题：

\begin{enumerate}
\item 能否出现 \(AC\perp BC\) 的情况？说明理由；
\item 证明过 \(A\)，\(B\)，\(C\) 三点的圆在 \(y\) 轴上截得的弦长为定值.
\end{enumerate}
\end{problem}

\begin{solution}
\begin{enumerate}
\item 设 \(A(x_1,0)\)，\(B(x_2,0)\).因为 \(x_1\)，\(x_2\) 是方程 \(x^2+mx-2=0\) 的两个根，由韦达定理得 \(x_1x_2=-2\).于是
\(\overrightarrow{AC}=(-x_1,1)\)，\(\overrightarrow{BC}=(-x_2,1)\)，从而
\(\overrightarrow{AC}\cdot\overrightarrow{BC}=x_1x_2+1=-2+1=-1\ne0\).所以不可能出现 \(AC\perp BC\) 的情况.
\item 设过 \(A\)，\(B\)，\(C\) 三点的圆的圆心为 \(O\left(-\dfrac m2,b\right)\)，半径为 \(r\).因为 \(AB\) 在 \(x\) 轴上，圆心的横坐标为 \(AB\) 中点的横坐标 \(\dfrac{x_1+x_2}{2}=-\dfrac m2\).又
\(r^2=b^2+\left(\dfrac{x_1-x_2}{2}\right)^2=b^2+\dfrac{(x_1+x_2)^2-4x_1x_2}{4}=b^2+\dfrac{m^2+8}{4}\).

点 \(C(0,1)\) 在圆上，所以
\(\dfrac{m^2}{4}+(1-b)^2=b^2+\dfrac{m^2+8}{4}\)，化简得 \(b=-\dfrac12\).因此圆的方程为
\(\left(x+\dfrac m2\right)^2+\left(y+\dfrac12\right)^2=\dfrac{m^2+9}{4}\).令 \(x=0\)，得 \(\left(y+\dfrac12\right)^2=\dfrac94\)，所以圆在 \(y\) 轴上的两个交点的纵坐标为 \(1\)，\(-2\)，弦长恒为 \(1-(-2)=3\).
\end{enumerate}
\end{solution}

\begin{problem}
已知函数 \(f(x)=\ln x+ax^2+(2a+1)x\).

\begin{enumerate}
\item 讨论 \(f(x)\) 的单调性；
\item 当 \(a<0\) 时，证明 \(f(x)\le -\dfrac{3}{4a}-2\).
\end{enumerate}
\end{problem}

\begin{solution}
\begin{enumerate}
\item 函数定义域为 \(x>0\)，且
\(f'(x)=\dfrac1x+2ax+2a+1=\dfrac{(2ax+1)(x+1)}{x}\).由于 \(x>0\)，\(x+1>0\)，所以 \(f'(x)\) 的符号由 \(2ax+1\) 决定.

当 \(a\geq0\) 时，\(2ax+1>0\)，故 \(f(x)\) 在 \((0,+\infty)\) 上单调递增.当 \(a<0\) 时，令 \(2ax+1=0\)，得 \(x_0=-\dfrac{1}{2a}>0\).在 \(0<x<x_0\) 时 \(f'(x)>0\)，在 \(x>x_0\) 时 \(f'(x)<0\)，所以 \(f(x)\) 在 \(\left(0,-\dfrac{1}{2a}\right)\) 上单调递增，在 \(\left(-\dfrac{1}{2a},+\infty\right)\) 上单调递减.
\item 当 \(a<0\) 时，由上问知 \(f(x)\) 的最大值为
\(f\left(-\dfrac{1}{2a}\right)=\ln\left(-\dfrac{1}{2a}\right)-1-\dfrac{1}{4a}\).令 \(t=-\dfrac{1}{2a}>0\)，则需证明
\(\ln t-1-\dfrac{1}{4a}\leq-\dfrac{3}{4a}-2\)，这等价于 \(\ln t+1-t\leq0\).

令 \(g(t)=\ln t+1-t\)，则 \(g'(t)=\dfrac1t-1=\dfrac{1-t}{t}\).因此 \(g(t)\) 在 \((0,1)\) 上递增，在 \((1,+\infty)\) 上递减，最大值为 \(g(1)=0\).所以 \(\ln t+1-t\leq0\)，从而
\(f(x)\leq f\left(-\dfrac{1}{2a}\right)\leq-\dfrac{3}{4a}-2\).
\end{enumerate}
\end{solution}

\section{选考题：解答题}

\begin{problem}
在直角坐标系 \(xOy\) 中，直线 \(l_1\) 的参数方程为
\examdisplaycases{}{
x=2+t,\\
y=kt
}
（\(t\) 为参数），直线 \(l_2\) 的参数方程为
\examdisplaycases{}{
x=-2+m,\\
y=\dfrac{m}{k}
}
（\(m\) 为参数）. 设 \(l_1\) 与 \(l_2\) 的交点为 \(P\)，当 \(k\) 变化时，\(P\) 的轨迹为曲线 \(C\).

\begin{enumerate}
\item 写出 \(C\) 的普通方程；
\item 以坐标原点为极点，\(x\) 轴正半轴为极轴建立极坐标系，设 \(l_3:\rho(\cos\theta+\sin\theta)-\sqrt2=0\)，\(M\) 为 \(l_3\) 与 \(C\) 的交点，求 \(M\) 的极径.
\end{enumerate}
\end{problem}

\begin{solution}
\begin{enumerate}
\item 由 \(l_1\) 的参数方程得 \(t=x-2\)，所以 \(y=k(x-2)\)；由 \(l_2\) 的参数方程得 \(m=x+2\)，所以 \(y=\dfrac{x+2}{k}\).因 \(k\ne0\)，交点坐标满足
\(y^2=(x-2)(x+2)=x^2-4\)，即 \(x^2-y^2=4\).同时 \(y\ne0\)，否则由两条参数方程会导致 \(x=2\) 或 \(x=-2\) 时参数关系不相容.因此 \(C\) 的普通方程为 \(x^2-y^2=4\)，并满足 \(y\ne0\)（等价地可写为 \(x\ne2\)，\(y\ne0\)）.
\item 在极坐标中 \(x=\rho\cos\theta\)，\(y=\rho\sin\theta\)，所以直线 \(l_3\) 化为 \(x+y=\sqrt2\).与 \(C\) 联立，得
\(x^2-y^2=(x+y)(x-y)=4\)，从而 \(x-y=2\sqrt2\).解得
\(x=\dfrac{3\sqrt2}{2}\)，\(y=-\dfrac{\sqrt2}{2}\).因此
\(\rho^2=x^2+y^2=\dfrac{18}{4}+\dfrac{2}{4}=5\)，极径取非负值，故 \(\rho=\sqrt5\).
\end{enumerate}
\end{solution}

\begin{problem}
已知函数 \(f(x)=|x+1|-|x-2|\).

\begin{enumerate}
\item 求不等式 \(f(x)\ge 1\) 的解集；
\item 若不等式 \(f(x)\ge x^2-x+m\) 的解集非空，求 \(m\) 的取值范围.
\end{enumerate}
\end{problem}

\begin{solution}
\begin{enumerate}
\item 按 \(x\) 与 \(-1\)，\(2\) 的位置分段，有
\(\displaystyle f(x)=
\begin{cases}
-3,&x\leq-1,\\
2x-1,&-1<x<2,\\
3,&x\geq2.
\end{cases}\)
当 \(x\leq-1\) 时，\(f(x)=-3<1\)；当 \(-1<x<2\) 时，\(2x-1\geq1\) 等价于 \(x\geq1\)；当 \(x\geq2\) 时，\(f(x)=3\geq1\).所以解集为 \([1,+\infty)\).
\item 原不等式有解，当且仅当存在 \(x\) 使 \(m\leq f(x)-x^2+x\).令 \(g(x)=f(x)-x^2+x\)，则
\(\displaystyle g(x)=
\begin{cases}
-x^2+x-3,&x\leq-1,\\
-x^2+3x-1,&-1<x<2,\\
-x^2+x+3,&x\geq2.
\end{cases}\)

在 \(x\leq-1\) 上，二次函数顶点在区间外且递增，最大值为 \(g(-1)=-5\).在 \(-1<x<2\) 上，顶点为 \(x=\dfrac32\)，最大值为 \(g\left(\dfrac32\right)=\dfrac54\).在 \(x\geq2\) 上函数递减，最大值为 \(g(2)=1\).故 \(g(x)\) 的全局最大值为 \(\dfrac54\)，原不等式有解等价于 \(m\leq\dfrac54\).因此 \(m\in\left(-\infty,\dfrac54\right]\).
\end{enumerate}
\end{solution}
